%vsgtu714



%  %%%

%%%   Самарский государственный технический университет

%%%   Кафедра прикладной математики и информатики

%%%

%%%   Преамбула для статьи в журнал "Вестник Самарского государственного

%%%   технического университета. Серия: «Физико-математические науки»"

%%%   Ориентирована под MikTEx 2.4 for Win32

%%%

%%%   Хотя сам журнал собирается под GNU/Linux на дистрибутиве TeXLive



\documentclass[11pt,a4paper,twoside]{article}



\usepackage[cp1251]{inputenc}

\usepackage[T2A]{fontenc}

\usepackage[english,russian]{babel}

\usepackage{samgtubib}

\usepackage[dvips]{graphicx}

\usepackage[multiple]{footmisc}



\usepackage{samgtu}

\renewcommand{\VestnikYear}{2009} % Год издания Вестника

\renewcommand{\VestnikNumber}{2\,(19)} % Номер Вестника

\renewcommand{\VestnikMonth}{Ноябрь} % Месяц выхода Вестника

\renewcommand{\VestnikData}{01.11.09} % Дата подписания Вестника в печать

\renewcommand{\VestnikUPL}{26,64} % Усл. печ. л.

\renewcommand{\VestnikUIL}{25,73} % Уч.-изд. л.

\renewcommand{\VestnikRegNumber}{479/09} % Регистрационный номер

\renewcommand{\VestnikOrder}{994} % Номер заказа







%

% \usepackage{textcomp}

%

%

% \begin{document}

%

%

%  \tableofcontents

%  \newpage



 \renewcommand{\firstpage}{260}

 \renewcommand{\lastpage}{263}

%

%

% %%%%%%%%%%%%%%%%%%%%%%%%%%%%%%%%%%%%%%%%%%%%%%%%%%%%%%%%%%%%%%%%%%%%%%%%%%%%%%%%%%%%

% %Информация о статье и об авторах на русском и английском языках



% Номер УДК

\udk{539.319}%Пластичность

\msc{74A10, 74C05}



\title{Остаточные напряжения в образцах из сплавов В95 и Д16Т после пневмодробеструйной обработки}% Полный заголовок

{Остаточные напряжения в образцах из сплавов В95 и Д16Т \dots}% Заголовок, который идёт в верхний колонтитул



\author{В.\,А.~Кирпич\"eв\inst{1}, Д.\,В.~Иванов\inst{1},  М.\,Н. Саушкин\inst{2}}% Автор(ы) для заголовка, если авторов несколько укать

{К\,и\,р\,п\,и\,ч\,\"e\,в~В.\,А., И\,в\,а\,н\,о\,в~Д.\,В., С\,а\,у\,ш\,к\,и\,н~М.\,Н.}% Автор(ы) для оглавления и колонтитула через \,



\institute{\sgau.

\and

\samgtu.}





\email{E-mails}{sopromat@ssau.ru, msaushkin@gmail.com}



\otherinf{\fullauthor{Виктор Алексеевич Кирпичёв}  \regalia{к.т.н., доц.} \position{докторант}

\dept{каф. сопротивления материалов}

\fullauthor{Денис Всеволодович Иванов} \position{аспирант} \dept{каф. сопротивления материалов}

\fullauthor{Михаил Николаевич Саушкин}  \regalia{к.ф.-м.н., доц.} \position{докторант} \dept{каф. прикладной математики и информатики}}





\otherenginf{\fullauthor{Viktor A. Kirpichev} \regalia{Ph. D. (Techn.)} \position{Doctoral Candidate}

\dept{Dept. of Resistance Materials}

\fullauthor{Denis V. Ivanov} \position{Postgraduate Student}

\dept{Dept. of Resistance Materials}

\fullauthor{Mikhail N. Saushkin} \regalia{Ph.\,D. (Phys. \& Math.)} \position{Doctoral Candidate}

\dept{Dept. of Applied Mathematics \& Computer Science}

}



\keywords{остаточные напряжения, концентраторы напряжений, пневмодробеструйная обработка, термоэкспозиция, релаксация остаточных напряжений}



\workabstract{Исследованы остаточные напряжения в гладких образцах и~образцах с~надрезом из сплавов {\rm В95} и~{\rm Д16Т} после пневмодробеструйной обработки, а также после пневмодробеструйной обработки и термоэкспозиции при температуре \mbox{$T=125${\rm\,\textcelsius}} в~течение $100$ часов. Установлено, что термоэкспозиция приводит к~существенно большей релаксации остаточных напряжений в~образцах из сплава {\rm В95}, чем из сплава {\rm Д16Т}.}



\engtitle{Residual Stresses in Samples of  D16T and V95  Alloy  after Air Shot Peening Treatment}



\engauthor{V.\,A.~Kirpichev\inst{1}, D.\,V.~Ivanov\inst{1}, M.\,N.~Saushkin\inst{2}}{K\,i\,r\,p\,i\,c\,h\,e\,v~V.\,A., I\,v\,a\,n\,o\,v~D.\,V., S\,a\,u\,s\,h\,k\,i\,n~M.\,N.}



\enginstitute{\engsgau.

\and

\engsamgtu.}







\engabstract{Residual stresses in smooth samples and samples with a cut cast out of {\rm V95} and {\rm D16T} alloys after air shot penning and also after air shot-penning and thermoexpositions are studied  at temperature $T=125\,^\circ $C during $100$ hours. It is established that thermoexposition leads to an essentially larger relaxation of residual stresses in samples cast out of  {\rm V95} alloys than those cast out of  {\rm D16T.}}



\engkeywords{residual stresses, Stress concentrators, air shot peening, thermoexposition, residual stresses relaxation}



\date{20/VII/2009} % Дата поступления статьи в редакцию.

\revisiondate{31/VIII/2009} % В окончательном виде



%%%%%%%%%%%%%%%%%%%%%%%%%%%%%%%%%%%%%%%%%%%%%%%%%%%%%%%%%%%%%%%%%%%%%%%%%%%%%%%%%%%%





\maketitle



\small





Исследовались остаточные напряжения в~образцах диаметром 15~мм с~отверстием

диаметром 5~мм из сплавов В95 и~Д16Т после упрочнения на пневмодробеструйной

установке при давлении воздуха~0,25~МПа стальными шариками 1,5--2,0~мм.

Механические характеристики исследуемых материалов приведены в~таблице.







\begin{table}[h!]

\centering \footnotesize

\textbf{Механические характеристики материалов}

\begin{tabular}{c|c|c|c|c|c}

\hline

Материал&

$\sigma _{0,2},$~МПа&

$\sigma _{\mbox{в}}$,~МПа&

$\delta$, \% &

$\Psi$, \% &

$S_{\text{к}}$,~МПа \\

\hline

В95&

645&

679&

6&

17,3&

755 \\

%\hline

Д16Т&

451&

618&

11,9&

13,3&

684 \\

\hline

\end{tabular} \vspace{-5mm}

\end{table}





\begin{wrapfigure}{O}{.55\textwidth} \centering %\vspace{-3mm}

\includegraphics[width=.5\textwidth]{kirp1.eps}

\caption{Распределение остаточных напряжений в~гладких образцах из сплавов

В95~(1, 3) и~Д16Т (2, 4): \scriptsize 1, 2 "--- после ПДО, 3, 4 "--- после ПДО и~термоэкспозиции}

\end{wrapfigure}



Осевые $\sigma _{z}$ и окружные $\sigma _{\theta }$ остаточные

напряжения в образцах определялись методом колец и~полосок [1, 2] как

непосредственно после пневмодробеструйной обработки (ПДО), так и после ПДО с~последующей термоэкспозицией в~течение $t=100$ часов при температуре $T=125$\,\textcelsius. Распределение остаточных напряжений в гладких образцах по толщине

поверхностного слоя $a$ приведено на рис.~1, из которого можно видеть, что в~гладких образцах после упрочнения

(кривые~1~и~2) остаточные напряжения для сплавов В95 и~Д16Т до глубины

$a=0,3$~мм различаются незначительно, при этом в~образцах из Д16Т сжимающие

остаточные напряжения несколько больше. Начиная с~глубины $a=0,3$~мм напряжения

в~образцах из сплава Д16Т выше, чем из~В95, причём смена знака остаточных

напряжений в образцах из сплава Д16Т происходит на существенно большей

глубине ($a=0,64$ мм) в~сравнении с~образцами из В95 ($a=0,44$~мм).











Б\'ольшая толщина слоя со сжимающими остаточными напряжениями в~уп\-роч\-нён\-ных

образцах из Д16Т объясняется, по"=видимому, б\'ольшей пластичностью этого

сплава по сравнению с~В95 (см.~таблицу), в~связи с~чем поверхностный слой

образцов из Д16Т деформируется при упрочнении на большую глубину.



Следовательно, после упрочнения на одинаковых режимах распределение

остаточных напряжений в образцах из сплава Д16Т является более полным, чем в~образцах из В95.

Это обстоятельство может оказать существенное влияние на сопротивление усталости.

В~соответствии с~закономерностями, установленными в~работе~[3],

приращение предела выносливости упрочнённых ПДО образцов из

сплава Д16Т будет выше, чем из сплава В95.







Термоэкспозиция упрочнённых ПДО образцов привела к снижению (релаксации)

сжимающих ос\-таточных напряжений, особенно в образцах из сплава В95. Если~в образцах из Д16Т максимальные напряжения до термоэкспозиции превышали

аналогичные напряжения в образцах из В95 на 8~МПа (340 и~332~МПа

соответственно), то после термоэкспозиции превышение составило 151 МПа (325~и~174~МПа).

Кроме того, в~образцах из сплава В95 глубина смены знака

остаточных напряжений уменьшилась на 0,08~мм, в то время как в~образцах из

Д16Т она практически не изменилась.

В~связи с~этим снижение характеристик

сопротивления усталости для образцов из сплава В95 после термоэкспозиции

будет больше, чем для образцов из Д16Т~[3].



Далее на упрочнённые ПДО цилиндрические образцы из тех же сплавов (В95 и~Д16Т)

таких же размеров (наружный диаметр "--- 15~мм, диаметр отверстия "--- 5~мм)

наносились круговые надрезы полукруглого профиля двух радиусов $R=0,3$ и~$R=0,5$~мм.

Надрезы наносились как на образцы непосредственно после

упрочнения, так и~после упрочнения с последующей термоэкспозицей в~течение

$t=100$ часов при температуре $T=125$\,\textcelsius.



Остаточные напряжения в~образцах с~надрезом определялись расчётным путём по

остаточным напряжениям гладких образцов (рис.~1) с~учётом перераспределения

остаточных усилий гладких образцов при нанесении надреза.

При этом использовались как аналитический~[4], так и~численный методы.

В~случае решения задачи численным методом применялся программный комплекс~{\tt ANSYS}.

Результаты расчётов распределения осевых $\sigma _{z}$ остаточных

напряжений в~образцах с~надрезом представлены на рис.~2 и~3, из которых

следует, что после нанесения надрезов на упрочнённые ПДО образцы на их дне возникают значительные по величине

сжимающие остаточные напряжения, которые для В95 составляют 686~МПа, а для Д16Т "---  720~МПа.

Следовательно, наибольшие напряжения выше не только предела

текучести, но и~предела прочности используемых в~исследовании материалов (см. таблицу).

Этому явлению дано объяснение в~работах [5, 6], где показано, что наибольшие остаточные

напряжения после поверхностного упрочнения могут быть выше сопротивления разрыву материала $S_{\text{к}}$ на 15{\%}.

В~нашем исследовании остаточные напряжения на дне надреза этого предела не превышают

(см.~таблицу).





\begin{figure}[b!]

\centering \vspace{-6mm}

 \includegraphics{kirp2-3.eps}

\parbox[t]{.48\textwidth}{\caption{Распределение остаточных напряжений в~образцах с~надрезом из сплава

 В95 после ПДО (1, 3), ПДО и термоэкспозиции (2, 4): \scriptsize 1, 2 "--- $R=0,3$~мм; 3, 4 "---

 $R=0,5$~мм}}\hfill

\parbox[t]{.48\textwidth}{\caption{Распределение остаточных напряжений в образцах с надрезом из сплава Д16Т после ПДО (1, 3) ПДО и~термоэкспозиции (2, 4): \scriptsize 1, 2 "--- $R=0,3$~мм; 3, 4 "---

 $R=0,5$~мм}} \vspace{-3mm}

\end{figure}





Во всех случаях остаточные напряжения в~образцах с~радиусом надреза $R=0,3$~мм

выше, чем с~$R=0,5$ мм, что объясняется, во-первых, большей концентрацией

напряжений и, во"=вторых, большими средними значениями остаточных напряжений

гладких образцов по толщине поверхностного слоя при $a=0,3$~мм, чем при $a=0,5$~мм

(см. рис.~1).

В~результате дополнительные остаточные напряжения, возникающие за

счёт перераспределения остаточных усилий гладкого образца при нанесении

надреза в~надрезанных образцах с~$R=0,3$~мм выше, чем с~$R=0,5$~мм.



Термоэкспозиция упрочнённых образцов привела к заметному снижению сжимающих

остаточных напряжений, особенно в образцах из сплава В95, так как в~гладких

упрочнённых образцах из этого сплава, на которые наносились надрезы,

остаточные напряжения значительно ниже, чем из сплава Д16Т.

Поэтому следует ожидать, что приращение предела выносливости для образцов из сплава В95 за

счёт ПДО будет ниже, чем для образцов из Д16Т~[3].



Таким образом, проведённое исследование показало, что при

пневмодробеструйной обработке на одних и тех же режимах как в~гладких

образцах, так и~в~образцах с~надрезом из сплава Д16Т толщина слоя со

сжимающими остаточными напряжениями существенно больше, чем из сплава В95.

Термоэкспозиция образцов в течение 100 часов при температуре $125$\,\textcelsius{}

приводит к~релаксации остаточных напряжений, причём действие термоэкспозиции

заметно выше для образцов из сплава В95.

Выявленные закономерности формирования остаточных напряжений в~результате упрочнения ПДО следует

учитывать при прогнозировании предела выносливости деталей из сплавов В95 и~Д16Т.



{\thanks{Работа выполнена при поддержке Федерального агентства по образованию $($проект РНП $2.1.1/3397)$ и РФФИ $($проект \rm \No~$07{-}01{-}00478{-}a).$}}



\vspace{3mm}



\begin{thebibliography}{9}

\RBibitem{kirp:1}

\by Иванов~С.\,И.

\paper К~определению остаточных напряжений в цилиндре методом колец и~полосок

\inbook Остаточные напряжения. {\rm Вып. 53}

\publaddr Куйбышев

\publ КуАИ

\yr 1971

\pages 32--42



\RBibitem{kirp:2}

\by Биргер~И.\,А.

\book Остаточные напряжения

\publaddr М.

\publ Машгиз

\yr 1963

\finalinfo 232~с





\RBibitem{kirp:3}

\by Павлов~В.\,Ф.

\paper Влияние характера распределения остаточных напряжений по толщине поверхностного слоя детали на сопротивление усталости

\jour Известия вузов. Машиностроение

\yr 1987

\issue 7

\pages 3--6



\RBibitem{kirp:4}

\by Иванов~С.\,И., Шатунов М.\,П., Павлов В.\,Ф.

\paper Влияние остаточных напряжений на выносливость образцов с~надрезом

\inbook Вопросы прочности элементов авиационных конструкций. {\rm Вып. 1}

\publaddr Куйбышев

\publ КуАИ

\yr 1974

\pages 88--95



\RBibitem{kirp:5}

\by  Радченко~В.\,П., Павлов~В.\,Ф.

\paper К~вопросу о наибольшей величине остаточных напряжений в~упрочнённых деталях

\inbook Проблемы и перспективы развития двигателестроения

\bookinfo Материалы докладов МНТК. Часть 1

\publaddr Самара

\publ СГАУ

\yr 2009

\pages 251--252



\RBibitem{kirp:6}

\by Павлов~В.\,Ф., Радченко~В.\,П.

\paper О наибольшей величине остаточных напряжений в~упрочнённых деталях

\serial Математическое моделирование и краевые задачи

\inbook Труды Шестой Всероссийской конференции с международным участием. {\rm Часть 1}

\bookinfo Математические модели механики, прочности и надёжности элементов конструкций

\publaddr Самара

\publ СамГТУ

\yr 2009

\pages 186--189

\end{thebibliography}





\makeengtitle





% \end{document}

